\chapter{Tool Runtime}
\label{ch:tools}

The tool runtime is the subsystem responsible for resolving, validating, and
invoking tool definitions at runbook execution time.  Tools are the primary
mechanism by which gert performs side-effecting work: executing binaries,
calling JSON-RPC services~\cite{jsonrpc-spec}, and delegating to MCP-hosted capabilities~\cite{mcp-spec}.

\section{Tool Discovery and Catalog Construction}
\label{sec:tool-discovery}

Tool resolution is split into two plan-time phases so that ``what tools
exist'' is decided before ``what a step means'' (AR-TP-2,
\texttt{.squad/decisions/archive/barbara-tool-packages-architecture-ruling.md}).

\begin{description}
  \item[Phase C --- Catalog construction] Runs once per run, before any step
        is planned. It produces the \texttt{ResolvedToolCatalog}: a map from
        fully-qualified tool name to a single \texttt{ToolDefinition}, plus a
        frozen package map (\S\ref{sec:tool-packages}). It is a pure function
        of the workspace files, \texttt{requires:} declarations, extension
        contributions, and MCP \texttt{tools/list} responses. It has no
        knowledge of any individual step.
  \item[Phase B --- Binding] Runs once per runbook file in the include
        closure. Each \texttt{toolRefs} entry in that file binds one catalog
        entry to one file-local name. Every \texttt{step.tool.name} in that
        file \textbf{MUST} resolve to a name bound in that same file
        (\texttt{PLAN-010}); tool-name binding is \textbf{lexically scoped} to
        the declaring runbook file and is never inherited from a parent that
        includes it (\S\ref{sec:tool-packages-includes}).
\end{description}

A runtime \textbf{MUST NOT} mutate the catalog after Phase C completes. Any
attempt to contribute a package to the catalog after the freeze --- including
a lazily-expanded include that would introduce a previously-undeclared
package --- is \texttt{PKG-017}.

\subsection{Tiers}
\label{sec:tool-tiers}

Tools enter the catalog from exactly five tiers. A higher tier number is
higher precedence.

\begin{center}
\begin{tabular}{cp{6.2cm}p{4.8cm}}
\hline
\textbf{Tier} & \textbf{Source} & \textbf{Namespacing} \\
\hline
0 & Built-in registry compiled into the host & bare names \\
1 & Packages from \texttt{requires:} (project-level, then runbook-level) & \texttt{<package>/<tool>} + bare name \\
2 & Project tools: \texttt{.gert/config.yaml} \texttt{tool-paths} (in listed order), then \texttt{<workspace>/tools/} & bare names \\
3 & Dynamic: extension contributions (\texttt{capability/tool-registration}), then MCP \texttt{tools/list} & \texttt{<extension>/<tool>}, \texttt{<server>/<tool>} --- mandatory \\
4 & \texttt{toolRefs[].path} in the runbook being planned & bare name, file-local \\
\hline
\end{tabular}
\end{center}

Every catalog entry carries \textbf{both} a fully-qualified name
(\texttt{<origin>/<tool>}) and, for tiers 0/1/2/4, a bare name. Tier 3
entries have \textbf{no} bare name --- dynamic contributions are always
qualified. This alone removes the entire class of ``extension silently
shadows kubectl'' failures.

Built-ins (tier 0) are always present; there is no ``disable builtins''
switch in the MVP. A package or project tool sharing a built-in's bare name
shadows it, and that shadowing is subject to the same declaration
requirement as any other cross-tier shadow (rule 2 below).

\subsection{Precedence Rule}
\label{sec:tool-precedence}

This replaces every prior ``first match wins'' or ``later entries take
precedence'' formulation in this specification. The direction is uniform
across the whole spec (\S\ref{sec:ext-discovery} adopts the identical rule):
\textbf{higher tier wins; ties are errors.}

\begin{enumerate}
  \item \textbf{Same-tier collision on the same bare name is a hard error}
        (\texttt{PKG-006}). There is no warning-and-continue. Two packages
        that both export the bare name \texttt{kubectl} collide; the runbook
        \textbf{MUST} disambiguate with \texttt{toolRefs[].package} or with
        the fully-qualified name.
  \item \textbf{Cross-tier shadowing is permitted, and silent, only when the
        shadowing entry is at a strictly higher tier AND the runbook
        declares the intent} via \texttt{toolRefs[].package} or
        \texttt{toolRefs[].path}. Undeclared cross-tier shadowing is
        \texttt{PKG-022} (a hard error, not a warning) --- this is the single
        intentionally breaking behaviour change introduced by this
        specification (\S\ref{sec:tool-packages-compat}).
  \item Fully-qualified names never collide and never shadow; they are
        always addressable even when the bare name is ambiguous.
\end{enumerate}

\subsection{Determinism of Enumeration}
\label{sec:tool-determinism}

Non-determinism from filesystem iteration order is a defect, not an
implementation detail. The following is normative:

\begin{itemize}
  \item Directory scans are recursive over \texttt{*.tool.yaml}, and results
        are sorted by the \textbf{workspace-relative POSIX path}, NFC-normalised,
        compared by Unicode code point, ascending.
  \item \texttt{requires:} entries are processed in declaration order;
        project-level \texttt{requires:} precede runbook-level
        \texttt{requires:}.
  \item \texttt{tool-paths} entries are processed in declaration order.
  \item Extension contributions are processed in extension-resolution order;
        within one extension, in the order returned by
        \texttt{contributions/list}.
  \item MCP servers are processed in declaration order; within one server,
        tools are sorted by name --- the server's array order is
        \textbf{not} trusted, since MCP servers make no ordering guarantee.
  \item Two runs over an identical file tree \textbf{MUST} produce
        byte-identical catalog digests (\S\ref{sec:tool-package-digest}).
\end{itemize}

\section{Runbook Tool Bindings: \texttt{requires:} and \texttt{toolRefs:}}
\label{sec:tool-bindings}

\subsection{\texttt{requires:} is canonical}

\texttt{requires:} is the \textbf{sole, canonical} mechanism by which a
project or a runbook declares a dependency on a tool package. It may appear
at the top level of \texttt{.gert/config.yaml} (project scope) and at the top
level of any runbook (runbook scope). A \texttt{requires:} entry
\textbf{supplies} tool definitions to the catalog (\S\ref{sec:tool-tiers},
tier 1).

A previously-discussed \texttt{toolPackages:} key was considered and
\textbf{rejected}: it duplicated the job of \texttt{requires:} under a
different name. \texttt{toolPackages:} \textbf{MUST NOT} be introduced in any
schema. Any document containing a \texttt{toolPackages:} key is rejected with
\texttt{PKG-020} (\texttt{unsupported key: use requires:}). There is
\textbf{no alias, no dual-read, and no deprecation window} for this key ---
it never shipped in any released schema, so the rejection is a permanent hard
error from day one, not a transitional one.

\subsection{\texttt{toolRefs:} binds and narrows}

\texttt{toolRefs:} declares which tool names a runbook file's \texttt{flow}
may reference, and optionally pins where each name is bound from. It does
\textbf{not} itself add packages to the catalog. The relationship is
one-way: \texttt{toolRefs} may \emph{narrow} what \texttt{requires:} supplied;
it may never widen it.

\begin{minted}{yaml}
requires:
  - package: acme.incident-tools
    version: "^1.4.0"
    path: ./vendor/acme-incident-tools   # optional here; project config may supply it

toolRefs:
  - name: kubectl
    package: acme.incident-tools    # pins tier and supplier; collision-proof by construction
    version: "^1.4.0"               # checked against the resolved tool's meta.version
\end{minted}

\subsubsection{\texttt{ToolRef} field disposition}

\begin{center}
\begin{tabular}{p{2cm}p{11.3cm}}
\hline
\textbf{Field} & \textbf{Disposition} \\
\hline
\texttt{name} & REQUIRED. The name used in \texttt{step.tool.name}. \\
\texttt{package} & \textbf{NEW.} Optional. Reverse-DNS package name that MUST have supplied this tool. Pins the tier and the supplier; makes a collision impossible by construction. Mutually exclusive with \texttt{path}. \\
\texttt{version} & \textbf{NEW.} Optional. Constraint (\S\ref{sec:tool-version-grammar}) checked against the resolved tool definition's \texttt{meta.version}. \\
\texttt{path} & RETAINED. Optional. Direct path to a \texttt{.tool.yaml}, resolved per \S\ref{sec:tool-secure-paths}. Highest tier (tier 4, \S\ref{sec:tool-tiers}). Mutually exclusive with \texttt{package}. \\
\texttt{actions} & RETAINED, \textbf{semantics changed}. Documentation and a plan-time reachability assertion only: every listed action MUST exist on the resolved tool (\texttt{PKG-012} if not). It is \textbf{NOT} an allowlist and \textbf{MUST NOT} gate invocation in the MVP. Emits deprecation warning \texttt{PKG-W001} when present. \\
\texttt{source} & RETAINED, \textbf{semantics changed}. Opaque provenance string. \textbf{MUST NOT} participate in resolution. Emits deprecation warning \texttt{PKG-W002}. Replaced by \texttt{package}. \\
\texttt{alias} & \textbf{REMOVED} from the schema. Out of scope for the MVP. Presence is \texttt{PKG-021}. \\
\hline
\end{tabular}
\end{center}

\paragraph{Deprecation schedule.} Per AR-TP-10, \texttt{source} (deprecation
D-001) and \texttt{actions} (deprecation D-002) remain schema-valid --- they
are not hard errors --- because fixtures \texttt{r01-k8s-incident} and
\texttt{r05-security-breach} use them today. They lose all resolution and
enforcement power immediately (\texttt{source} never participated in
resolution beyond documentation; \texttt{actions} is now non-enforcing).
Removal is scheduled for the first release after the fixture corpus is
migrated to \texttt{package}. This is a numbered deprecation with a removal
target, not an open-ended ``eventually''.

\subsection{Duplicate and conflicting declarations}

\begin{itemize}
  \item The same \texttt{package} name appearing twice within one
        \texttt{requires:} list is \texttt{PKG-005}, unconditionally --- even
        if the constraints and paths are identical. Duplicates are always an
        authoring mistake.
  \item The same \texttt{package} name in the project-level and
        runbook-level \texttt{requires:} lists is legal; the constraints are
        \textbf{intersected}. An empty intersection is \texttt{PKG-002}.
  \item A runbook-level \texttt{requires:} entry with no \texttt{path} and no
        matching project-level entry is \texttt{PKG-001}.
  \item A runbook-level entry whose \texttt{path} disagrees with the
        project-level entry for the same package name is \texttt{PKG-024}
        (conflicting binding) --- the project is authoritative and
        disagreement is an error, not a silent override.
\end{itemize}

\section{Name Resolution Order (Phase B)}
\label{sec:tool-name-resolution}

When a runbook step references a tool by name, the host resolves the bare (or
qualified) name against the frozen catalog produced by Phase C, using the
tier precedence and collision rules of \S\ref{sec:tool-precedence}. The
resolution is a lookup, not a search: because Phase C already applied the
five-tier precedence and rejected same-tier collisions and undeclared
cross-tier shadows, Phase B never needs its own precedence logic --- a bare
name in the frozen catalog is unambiguous by construction, or the run failed
to plan.

Tool authors SHOULD still use reverse-DNS namespace prefixes for their
package names (e.g., \texttt{acme.pd-lookup}) to make qualified references
self-describing, even though the tier/collision rules no longer depend on
naming convention for correctness.

\section{Tool Definition Schema}
\label{sec:tool-definition-schema}

\texttt{actions:} is \textbf{always an array} of entries keyed by
\texttt{name:} --- never a mapping keyed by the action name. This is the one
canonical \texttt{tool/v1} action representation: the same shape is used
whether the action is process-backed or substituted by a runbook
(\S\ref{sec:tool-substitution}), so a reader never has to reconcile two
competing normative shapes for the same key.

\begin{minted}{yaml}
# tools/my-service.tool.yaml
apiVersion: tool/v1

meta:
  name: my-service              # Unique identifier within the workspace
  version: "2.0.0"             # Tool definition version (semver)
  description: "Calls the MyService REST API"
  binary: my-service-cli        # Executable name (looked up in PATH)

transport:
  mode: stdio-jsonrpc           # stdio | stdio-jsonrpc | grpc | mcp

  # For stdio-jsonrpc and mcp: startup behaviour
  startup:
    argv: ["--server-mode"]     # Extra args passed to the binary
    ready-pattern: "ready"      # Regex matched against stdout/stderr to
                                # confirm the server is accepting requests
    timeout: "10s"
    shutdown-method: "shutdown" # JSON-RPC method to call on graceful stop

governance:
  requires-capabilities: []     # Host capabilities the tool needs, e.g.
                                # ["capability/network"]
  allowed-environments: ["real", "dry-run"]
  requires-approval: false      # Set true to gate invocation on operator OK

# Optional: per-platform implementation blocks (mobile execution only).
# This is TOOL-LEVEL impl: --- platform dispatch. It is a distinct key from
# the ACTION-LEVEL execute: block below (§Action Substitution); see the
# disambiguation note at the end of this section for how the two interact.
impl:
  ios:
    transport: native-sdk       # Dispatch to registered iOS handler
    handler: GertSDK.MyService.lookup
  android:
    transport: native-sdk       # Dispatch to registered Android handler
    handler: com.gert.tools.myservice.LookupHandler

actions:
  - name: lookup                # REQUIRED. Unique within this tool.
    description: "Look up an entity by ID"
    argv: ["lookup", "--id", "{{ .id }}"]  # Only for stdio transport;
                                            # rendered from args: below
    timeout: "30s"
    args:                        # Typed argument/input declarations. The
                                 # SINGLE vocabulary for what step.tool.args
                                 # may contain, for both process and
                                 # substituted actions (see below).
      id:
        type: string
        required: true
        description: "Entity identifier"
        redact: false
        enum: ["prod-id", "stage-id"]  # OPTIONAL value constraint, string only
    output:                     # Raw output shape. Process actions only:
      format: json               # text | json | jsonl. Describes the shape
                                  # of the captured stdout stream, not a set
                                  # of named values (see 'outputs:' below).
    capture:
      stdout:
        format: json
\end{minted}

\paragraph{\texttt{args:} vs \texttt{inputs:}, \texttt{output:} vs
\texttt{outputs:} --- one vocabulary, stated normatively.} Every action,
regardless of \texttt{execute.kind} (\S\ref{sec:tool-substitution}),
declares its callable parameters under \texttt{args:} --- a map from
argument name to \texttt{\{type, required, description, redact, enum\}},
exactly as shown above. \texttt{step.tool.args} in a calling runbook \textbf{ALWAYS}
binds by name to this same \texttt{args:} map, whether the resolved action
is process-backed or substituted; there is no separate \texttt{inputs:}
vocabulary and no kind-dependent binding rule. (An earlier draft of the
substitution contract used \texttt{inputs:} for this purpose on substituted
actions only, which left the binding target of \texttt{step.tool.args}
unstated for that case; that draft vocabulary is superseded by this
section and does not appear elsewhere in this specification.)

\paragraph{\texttt{enum} on \texttt{args.<name>} and \texttt{outputs.<name>}.}
Either key MAY declare \texttt{enum:}, valid only when the declaration
resolves to \texttt{type: string}; these are two of the four sites at which
\texttt{enum} is normatively defined across the specification
(\texttt{design/gert/sections/03-schema-vnext.tex} \S\ref{para:enum-constraint}
states the full well-formedness, Unicode, ordering, default, plan/runtime
timing, replay, and metadata/redaction rules, which apply identically here;
the other two sites are the runbook \texttt{inputs.<name>}/\texttt{outputs.<name>}
declarations in that same section). \texttt{enum} on \texttt{args.<name>} is
checked at the same argument-binding sites as any other typed
\texttt{args:} entry (\S\ref{sec:tool-substitution} \S Input contract);
\texttt{enum} on \texttt{outputs.<name>} is checked at production time, once
the substitute has produced its own \texttt{outputs.<name>.value}
(\S\ref{sec:tool-substitution-outputs} below). \texttt{enum} on
\texttt{type: secret} is forbidden (\texttt{ENUM-001}); \texttt{args:}
entries marked \texttt{redact: true} MAY declare \texttt{enum}, but its
member list is redacted in every trace, JSON, and diagnostic surface
(\texttt{design/gert/sections/08-security-and-trust.tex} \S Secret
Resolution and Redaction).

The two output-shaped keys are \textbf{not} the same vocabulary, because
they describe different things and both are needed:

\begin{itemize}
  \item \texttt{output:} (singular) --- \textbf{process actions only}.
        Describes the raw shape of the captured stdout stream
        (\texttt{text}\,|\,\texttt{json}\,|\,\texttt{jsonl}), consumed by the
        \texttt{capture:} block's GCP paths (\texttt{stdout},
        \texttt{json.*}, etc., \S\ref{sec:tool-substitution-outputs}).
  \item \texttt{outputs:} (plural) --- \textbf{required for substituted
        (\texttt{execute.kind: runbook}) actions}; not applicable to process
        actions. Declares the named, typed values the substitute
        \textbf{MUST} produce (\S\ref{sec:tool-substitution-outputs}), since
        a substitute produces a set of named values rather than a single
        captured stream.
\end{itemize}

A tool definition schema (hand-authored, consistent with the other
\texttt{design/gert/schemas/} artefacts) MAY be added in a future revision;
until then this section and \S\ref{sec:tool-substitution} are the sole
normative source for \texttt{.tool.yaml} structure, and both use the shapes
above exclusively.

\subsection{Transport-Specific Fields}

\paragraph{stdio} The binary is spawned once per invocation.  The
\texttt{argv} array in the action definition is rendered as a Go template with
the resolved argument values, then appended to the binary name to form the
complete command.  The process's stdout is captured as the tool output.

\paragraph{stdio-jsonrpc} The binary is spawned once per run (persistent
process).  The host keeps a single process alive across all action calls for
that tool.  Each action invocation sends a JSON-RPC request to the process's
stdin and reads the response from stdout.  The \texttt{startup.ready-pattern}
confirms the server is accepting connections.

\paragraph{mcp} The host spawns the binary and performs the MCP initialization
handshake (Section~\ref{sec:mcp-tools}).  Actions map to MCP \texttt{tools/call}
requests.

\paragraph{mcp-http} The host connects to a remote MCP server over HTTPS and
performs the Streamable HTTP initialization handshake
(Section~\ref{sec:mcp-http}).  Actions map to MCP \texttt{tools/call} HTTP
requests.  No local binary is spawned.  The \texttt{url:} field is required and
must use \texttt{https://} (plain HTTP is rejected).  Optional \texttt{auth:}
configures bearer-token acquisition; see Section~\ref{sec:mcp-http-auth}.

\section{Transport Layer}
\label{sec:tool-transport}

\subsection{stdio Transport}

\paragraph{Process establishment.}
The host calls \texttt{exec.Command(binary, resolvedArgv...)} with a
controlled environment.  A per-invocation timeout context wraps the call.

\paragraph{Invocation envelope.}
For stdio tools the invocation is entirely encoded in the process arguments
and environment; there is no envelope.  The arguments are rendered from the
action's \texttt{argv} template.

\paragraph{Response envelope.}
Stdout is captured in full after process exit.  The exit code is checked:
non-zero is mapped to a tool error.

\paragraph{Error envelope.}
\begin{minted}{json}
{
  "error": {
    "code": "TOOL_NONZERO_EXIT",
    "message": "tool exited with code 1",
    "details": {
      "exitCode": 1,
      "stderr": "connection refused\n",
      "durationMs": 320
    }
  }
}
\end{minted}

\subsection{stdio-jsonrpc Transport}

\paragraph{Process establishment.}
A single process is spawned at first use and kept alive for the duration of the
run.  The host waits for \texttt{startup.ready-pattern} on stdout/stderr before
sending the first request.

\paragraph{Invocation envelope.}
\begin{minted}{json}
{
  "jsonrpc": "2.0",
  "id": "<invocationId>",
  "method": "<actionName>",
  "params": {
    "args": { "id": "INC-12345" },
    "context": {
      "runId": "run-abc123",
      "stepId": "lookup_incident",
      "traceId": "trace-xyz",
      "userId": "ops@acme.example",
      "attempt": 1
    }
  }
}
\end{minted}

\paragraph{Response envelope.}
\begin{minted}{json}
{
  "jsonrpc": "2.0",
  "id": "<invocationId>",
  "result": {
    "output": { ... },
    "telemetry": {
      "startedAt": "2026-04-18T10:00:00Z",
      "finishedAt": "2026-04-18T10:00:00.320Z",
      "durationMs": 320
    }
  }
}
\end{minted}

\paragraph{Error envelope.}
\begin{minted}{json}
{
  "jsonrpc": "2.0",
  "id": "<invocationId>",
  "error": {
    "code": -32050,
    "message": "entity not found",
    "data": {
      "category": "runtime",
      "retryable": false,
      "details": { "entityId": "INC-12345" }
    }
  }
}
\end{minted}

Error code ranges:

\begin{itemize}
  \item \texttt{-32700} to \texttt{-32600}: JSON-RPC parse and request errors.
  \item \texttt{-32050}: Application-level tool error (retryable per
        \texttt{data.retryable}).
  \item \texttt{-32051}: Capability not granted.
  \item \texttt{-32052}: Input validation failed.
  \item \texttt{-32053}: Output schema validation failed.
  \item \texttt{-32054}: Timeout.
  \item \texttt{-32055}: Cancelled by host.
  \item \texttt{-32099}: Tool process crashed.
\end{itemize}

\subsection{MCP Transport}

See Section~\ref{sec:mcp-tools} for the full MCP integration specification.

\subsection{MCP Streamable HTTP Transport}
\label{sec:mcp-http}

The \texttt{mcp-http} transport connects to a remote MCP server over HTTPS
using the Streamable HTTP transport defined in MCP specification version
\texttt{2025-03-26}.  Unlike \texttt{mode: mcp} (which spawns a local
subprocess), \texttt{mcp-http} makes no assumption about where the server
runs or how it was deployed.

\subsubsection{Transport Fields}

\begin{minted}{yaml}
transport:
  mode: mcp-http
  url: https://icm-mcp-prod.azure-api.net/v1/
  auth:                             # optional; omit for unauthenticated servers
    provider: azure-cli
    scope: api://icmmcpapi-prod/mcp.tools
    allowed_hosts:
      - icm-mcp-prod.azure-api.net  # see Section~\ref{sec:mcp-http-auth}
\end{minted}

\begin{itemize}
  \item \textbf{\texttt{url:}} Required.  Must use \texttt{https://}; plain
        HTTP is rejected at validation time with error \texttt{MCP-001}.
        There is no escape hatch for non-TLS endpoints.
  \item \textbf{\texttt{auth:}} Optional.  When absent, no \texttt{Authorization}
        header is sent and \texttt{allowed\_hosts} is not required.  When
        present, see Section~\ref{sec:mcp-http-auth}.
\end{itemize}

Fields belonging to the \texttt{mcp} (stdio) arm --- \texttt{command:},
\texttt{args:} --- are rejected (not silently ignored) when
\texttt{mode: mcp-http} is declared.  This follows the same principle as
\texttt{expand} on a static include: a field the author clearly intended to
use for a different mode is an authoring error, not a no-op.

\subsubsection{Protocol Version}
\label{sec:mcp-http-protocol-version}

\texttt{mode: mcp-http} uses MCP protocol version \texttt{2025-03-26} (the
Streamable HTTP revision) via the \texttt{MCP-Protocol-Version} request header.
\texttt{mode: mcp} (stdio) uses version \texttt{2024-11-05} (the stdio framing
model) and the two are \emph{deliberately different}: they are different
transport specifications, not two implementations of the same protocol at
different revisions.  A server implementing one does not necessarily implement
the other, and changing either version risks breaking existing integrations.

\subsubsection{Redirect Behaviour}
\label{sec:mcp-http-redirects}

Redirect-following is \textbf{disabled} for authenticated \texttt{mcp-http}
requests (\texttt{http.ErrUseLastResponse}).  If a server responds with a
redirect, the request fails with \texttt{MCP-013}.  The remedy is to set
\texttt{url:} to the final endpoint; do not rely on the server to redirect.
Unauthenticated \texttt{mcp-http} requests may follow redirects normally.

The motivation: a redirect to an attacker-controlled host would bypass the
\texttt{allowed\_hosts} check and deliver a bearer token to an unapproved
destination.  Disabling redirects closes that path without requiring
additional runtime validation.

\subsubsection{Authentication}
\label{sec:mcp-http-auth}

See Section~\ref{sec:mcp-http-auth-detail} for the full authentication
specification.

\section{Invocation Context}
\label{sec:tool-context}

Every tool invocation — regardless of transport — receives a structured
context object.  For \texttt{stdio} tools, relevant fields are injected as
environment variables prefixed \texttt{GERT\_}; for \texttt{stdio-jsonrpc} and
\texttt{mcp} tools, the context is transmitted in the invocation envelope.

\begin{minted}{json}
{
  "context": {
    "runId":       "run-abc123",        // Globally unique run identifier
    "stepId":      "lookup_incident",   // Step ID within the runbook
    "traceId":     "trace-xyz789",      // Distributed trace correlation ID
    "userId":      "ops@acme.example",  // Identity executing the run
    "mode":        "real",              // real | dry-run | replay
    "attempt":     1,                   // Retry attempt number (1-based)
    "capabilities": [                   // Effective host capability set
      "capability/network",
      "capability/file-read"
    ]
  }
}
\end{minted}

\section{Capability Gate}
\label{sec:tool-capability-gate}

Before dispatching any tool invocation the host checks two conditions:

\begin{enumerate}
  \item The tool's \texttt{governance.requires-capabilities} list is a subset
        of the current run's granted capability set.  If not, the step fails
        immediately with a \texttt{CAPABILITY\_DENIED} error; the run does not
        attempt to start the tool process.
  \item The tool's \texttt{governance.allowed-environments} list includes the
        current execution mode.  A tool not listed for \texttt{dry-run} will be
        skipped (with a warning) when the run mode is dry-run.
\end{enumerate}

\section{Output Capture}
\label{sec:tool-output}

Tool output is mapped to named step captures using the \texttt{capture:}
declaration in the action definition.

\begin{minted}{yaml}
capture:
  dns_output: stdout          # Capture stdout as variable dns_output
  error_log:  stderr          # Capture stderr as variable error_log
  exit_code:  exitCode        # Capture numeric exit code
\end{minted}

For \texttt{json} output format, capture paths may use dot notation to extract
nested fields:

\begin{minted}{yaml}
capture:
  incident_id: "stdout.incident.id"
  severity:    "stdout.incident.severity"
\end{minted}

\paragraph{Redaction.}
Before any captured value is stored in the run state or emitted to the trace,
the governance redaction pipeline is applied.  Redaction rules are declared in
the runbook's \texttt{meta.governance.redact} block and in the tool action's
per-argument \texttt{redact: true} flag.  Redacted values are replaced with
\texttt{<redacted>} in all stored and displayed output.

\section{MCP Tool Integration}
\label{sec:mcp-tools}

MCP (Model Context Protocol) servers are long-lived processes that expose a
dynamic catalog of tools.  Gert treats an MCP server as a tool source, not as
a single tool.

\subsection{MCP Server Declaration}

\begin{minted}{yaml}
# In runbook meta or .gert/config.yaml
tools:
  mcp-servers:
    - name: acme-mcp
      binary: acme-mcp-server
      args: ["--port", "stdio"]
      transport: mcp
      startup:
        timeout: "15s"
\end{minted}

\subsection{MCP Discovery}

On run startup the host:

\begin{enumerate}
  \item Spawns the MCP server binary.
  \item Sends the MCP \texttt{initialize} request (per MCP specification).
  \item After receiving \texttt{initialized}, sends \texttt{tools/list}.
  \item Registers each returned tool under the name
        \texttt{<server-name>/<tool-name>} in the dynamic tool catalog.
\end{enumerate}

\subsection{MCP Tool Invocation}

When a runbook step calls an MCP-backed tool, the host sends:

\begin{minted}{json}
{
  "jsonrpc": "2.0",
  "id": "5",
  "method": "tools/call",
  "params": {
    "name": "lookup_incident",
    "arguments": { "id": "INC-12345" },
    "_gert_context": {
      "runId": "run-abc123",
      "stepId": "triage_step",
      "traceId": "trace-xyz"
    }
  }
}
\end{minted}

The \texttt{\_gert\_context} field is a gert extension to the MCP envelope.
MCP servers that do not recognise it MUST ignore it.  The response follows the
standard MCP \texttt{tools/call} response format.

\subsection{MCP vs stdio-jsonrpc Differences}

\begin{itemize}
  \item MCP servers expose multiple tools under a single process; stdio-jsonrpc
        tools are one tool per process.
  \item MCP uses a distinct initialization handshake (\texttt{initialize} /
        \texttt{initialized}) before tool registration.
  \item MCP tool names are discovered dynamically at runtime; stdio-jsonrpc
        action names are declared statically in the \texttt{.tool.yaml}.
  \item MCP responses use a \texttt{content} array; gert flattens text content
        to a single string for capture.
  \item \textbf{Protocol version:} \texttt{mode: mcp} (stdio) uses
        \texttt{MCP-Protocol-Version: 2024-11-05}; \texttt{mode: mcp-http}
        uses \texttt{MCP-Protocol-Version: 2025-03-26}.  This is intentional
        --- the two transport specifications are distinct, not two versions of
        the same protocol.  Changing either version breaks compatibility with
        servers built against that transport.
\end{itemize}

\subsection{MCP Streamable HTTP Authentication}
\label{sec:mcp-http-auth-detail}

\subsubsection{Overview}

The \texttt{auth:} block in a \texttt{mcp-http} transport declaration names a
\emph{provider} and an OAuth2 \emph{scope}.  It never contains a credential,
token, or secret --- by design.  The provider name selects the token-acquisition
mechanism (e.g.\ \texttt{azure-cli}); the scope identifies the target resource.
Gert acquires the token at runtime, attaches it as
\texttt{Authorization: Bearer \ldots} on each request, and never records the
token value in any log, trace event, or evidence store.

\begin{minted}{yaml}
auth:
  provider: azure-cli
  scope: api://icmmcpapi-prod/mcp.tools
  allowed_hosts:
    - icm-mcp-prod.azure-api.net
\end{minted}

\subsubsection{\texttt{allowed\_hosts} --- What It Protects Against}
\label{sec:mcp-http-allowed-hosts}

\texttt{allowed\_hosts} is \textbf{required} whenever \texttt{auth:} is
configured (\texttt{MCP-010}).  An empty or omitted list is a validation
error.

\paragraph{Why audience-scoped tokens are not sufficient alone.}
A bearer token acquired with \texttt{scope: api://icmmcpapi-prod/mcp.tools}
carries that audience in its claims.  A properly configured Azure AD resource
server that is \emph{not} the intended audience rejects the token.  This
prevents an attacker from \emph{consuming} the token at their own service.

However, it does \emph{not} prevent an attacker from \emph{replaying} the token
against the real, legitimate endpoint.  Possession is authorization: a token
sent to an attacker-controlled host can be exfiltrated and then presented to
the real IcM endpoint, impersonating the operator for the token's full
lifetime.  This is the confused-deputy / token-replay pattern.

\texttt{allowed\_hosts} is the prevention mechanism: the HTTP transport checks
the request URL's hostname against the declared list before attaching the
\texttt{Authorization} header.  If the hostname is not listed, the request
fails with \texttt{MCP-012} rather than delivering the token to an unapproved
host.

\paragraph{Matching semantics.}
Matching is \textbf{exact, case-insensitive, on the parsed hostname with port
stripped}.  Wildcards (e.g.\ \texttt{*.azure-api.net}) are not supported.  A
substring or prefix match is explicitly not used: if the allowed entry is
\texttt{icm-mcp-prod.azure-api.net}, then
\texttt{icm-mcp-prod.azure-api.net.evil.com} does \emph{not} match.  Authors
who need to permit multiple hosts must list each one explicitly.

The same semantics apply at static validation time (parse-time \texttt{MCP-011}
check on \texttt{url:}'s hostname) and at runtime (the \texttt{MCP-012} check
in \texttt{TokenGate.AttachToken} before each request).  The two checks use
identical logic so a config that passes validation cannot fail at runtime due
to a matching discrepancy.

\paragraph{Listing multiple hosts.}
\begin{minted}{yaml}
allowed_hosts:
  - icm-mcp-prod.azure-api.net
  - icm-mcp-staging.azure-api.net
\end{minted}

List only the hosts that will legitimately receive a bearer token.  A host
that should never receive the token --- a proxy, a health-check endpoint, a
CDN origin --- should not be listed; route it through an unauthenticated
transport instead.

\subsubsection{Unauthenticated Transports}

When \texttt{auth:} is absent, no \texttt{Authorization} header is sent and
\texttt{allowed\_hosts} is neither required nor checked.  Use the
unauthenticated form for MCP servers that use API-key headers, client
certificates, or network-level auth, as long as no bearer token is involved.

\subsubsection{Audit Trace}

Before each authenticated request, gert emits an \texttt{mcp/authAttached}
trace event containing \texttt{url\_host} and \texttt{scope} only --- never
the token value.  This enables post-hoc audit detection of scope-vs-host
anomalies without storing credentials.

\subsubsection{Error Codes}

\begin{description}
  \item[\texttt{MCP-001}] \texttt{url:} absent or not using \texttt{https://}.
  \item[\texttt{MCP-002}] \texttt{auth.provider} is not a recognized provider name.
  \item[\texttt{MCP-007}] Token acquisition failed (provider error, not logged-in, no consent).
  \item[\texttt{MCP-010}] \texttt{auth:} is configured but \texttt{allowed\_hosts} is absent.
  \item[\texttt{MCP-011}] Static check: \texttt{url:}'s hostname is not in \texttt{allowed\_hosts} (caught at parse time).
  \item[\texttt{MCP-012}] Runtime check: request hostname not in \texttt{allowed\_hosts} (closes the redirect path).
  \item[\texttt{MCP-013}] Authenticated request received a redirect response --- redirects are disabled on authenticated requests; update \texttt{url:} to the final endpoint.
\end{description}

\section{Cancellation and Timeout}
\label{sec:tool-cancellation}

\paragraph{Per-tool timeout.}
Declared in the action definition under \texttt{timeout:} (e.g., \texttt{30s}).
The default is the runbook's \texttt{meta.defaults.timeout}.  The host wraps
each invocation in a \texttt{context.WithTimeout}; when the deadline elapses
the host sends a \texttt{tools/cancel} request (JSON-RPC transport) or
sends SIGTERM (stdio transport), then forcibly terminates the process after a
5-second grace period.

\paragraph{Cancellation propagation.}
When a run is cancelled (via \texttt{Ctrl-C}, the debugger \texttt{quit}
command, or an API cancellation request), the host:

\begin{enumerate}
  \item Sets the run-level cancel context.
  \item Sends \texttt{tools/cancel} to all in-flight JSON-RPC and MCP
        invocations with a \texttt{reason: "run-cancelled"} field.
  \item After a 5-second drain window, forcibly terminates any unresponsive
        tool processes.
  \item Emits a \texttt{run.cancelled} trace event recording which steps were
        interrupted.
\end{enumerate}

The \texttt{tools/cancel} request envelope:

\begin{minted}{json}
{
  "jsonrpc": "2.0",
  "id": "cancel-5",
  "method": "tools/cancel",
  "params": {
    "invocationId": "5",
    "reason": "run-cancelled"
  }
}
\end{minted}

\section{Version Constraint Grammar}
\label{sec:tool-version-grammar}

\subsection{Versions}

Package and tool versions \textbf{MUST} be strict SemVer 2.0.0. A
non-conforming version is \texttt{PKG-025}. A leading \texttt{v} is
\textbf{rejected}, not stripped --- \texttt{v1.4.2} is not a valid version.

\subsection{Constraint grammar}

The following EBNF is the normative constraint grammar for the MVP. It is
deliberately \textbf{not} an expression language, and is therefore documented
here rather than as a \texttt{.ebnf} grammar file alongside GXL/GIS/GCP.

\begin{minted}{text}
Constraint   = Comparator , { " " , Comparator } ;
Comparator   = Caret | Tilde | Relational | Exact ;
Caret        = "^" , Version ;
Tilde        = "~" , Version ;
Relational   = ( ">=" | "<=" | ">" | "<" ) , Version ;
Exact        = [ "=" ] , Version ;
Version      = Major , "." , Minor , "." , Patch , [ "-" , Prerelease ] ;
\end{minted}

Space-separated comparators are \textbf{conjunctive (AND)}. There is no
disjunction in the MVP grammar.

\textbf{Not supported} (each is \texttt{PKG-003}): \texttt{||}, \texttt{*},
\texttt{x}/\texttt{X} wildcards, hyphen ranges (\texttt{1.2.3 - 1.5.0}),
partial versions (\texttt{\^{}1}, \texttt{\textasciitilde{}1.2}),
\texttt{!=}, whitespace variants other than a single ASCII space, and build
metadata inside a constraint. Partial versions are excluded deliberately:
\texttt{\^{}1} versus \texttt{\^{}1.0.0} differ across ecosystems, and
ambiguity in a governance-bearing artefact is unacceptable.

\subsection{Semantics}

\begin{itemize}
  \item \texttt{\^{}1.4.2} $\to$ \texttt{>=1.4.2 <2.0.0}. For \texttt{0.y.z}:
        \texttt{\^{}0.4.2} $\to$ \texttt{>=0.4.2 <0.5.0};
        \texttt{\^{}0.0.3} $\to$ \texttt{>=0.0.3 <0.0.4}. The \texttt{0.x}
        narrowing is stated explicitly because it is the most commonly
        mis-implemented rule across ecosystems.
  \item \texttt{\textasciitilde{}1.4.2} $\to$ \texttt{>=1.4.2 <1.5.0}.
  \item Ordering is SemVer 2.0.0 \S11 precedence.
  \item \textbf{Build metadata} (\texttt{+meta}) is ignored for ordering and
        for satisfaction, and \textbf{MUST NOT} be used for selection. It is
        retained verbatim in the lock and trace for provenance.
  \item \textbf{Prereleases:} a version with a prerelease component satisfies
        a comparator only if the comparator's own version has a prerelease
        component \textbf{and} shares the same \texttt{major.minor.patch}.
        So \texttt{1.5.0-rc.1} does \textbf{not} satisfy \texttt{\^{}1.4.0},
        but does satisfy \texttt{>=1.5.0-rc.1 <1.6.0}. This is the
        Cargo/npm rule and is the safe one for operational runbooks.
\end{itemize}

\subsection{Constraint checking, not solving}

The MVP has exactly one candidate per package: the local path. The
constraint is a \textbf{predicate on that candidate}, evaluated after loading
the manifest. There is no backtracking, no version selection, and no SAT
solving. Failure is \texttt{PKG-002} and \textbf{MUST} report the package
name, the resolved version, the effective (intersected) constraint, and the
list of declaration sites that contributed to the intersection.

% ═════════════════════════════════════════════════════════════════════════
\section{Tool Packages}
\label{sec:tool-packages}
% ═════════════════════════════════════════════════════════════════════════

A \textbf{tool package} is a local, on-disk directory that bundles one or
more \texttt{.tool.yaml} definitions behind an explicit export list. Tool
packages are the MVP's only supported distribution mechanism for tools ---
remote catalogs, registries, publishing workflows, transitive package
dependencies, and packaging of non-tool assets (docs, dashboards, policy
bundles, binaries-as-payload) are explicit non-goals of this MVP
(\S\ref{sec:tool-packages-nongoals}).

All new schemas in this section are hand-authored, language-neutral JSON
Schema (Draft 2020-12), consistent with \texttt{runbook.v1.schema.json},
with \texttt{additionalProperties: false} throughout.

\subsection{Package manifest --- \texttt{gert-package.yaml}}
\label{sec:tool-package-manifest}

Schema: \texttt{design/gert/schemas/tool-package.v1.schema.json}
(\texttt{apiVersion: tool-package/v1}).

\begin{minted}{yaml}
apiVersion: tool-package/v1     # REQUIRED, const "tool-package/v1"
meta:                           # REQUIRED
  name: acme.incident-tools     # REQUIRED. Reverse-DNS:
                                 # ^[a-z0-9]([a-z0-9-]*[a-z0-9])?
                                 #   (\.[a-z0-9]([a-z0-9-]*[a-z0-9])?)+$
  version: "1.4.2"              # REQUIRED. Strict SemVer 2.0.0
  description: "..."            # OPTIONAL
  license: "Apache-2.0"         # OPTIONAL, SPDX identifier
exports:                        # REQUIRED
  tools:                        # REQUIRED, minItems 1, unique by id and path
    - id: kubectl               # REQUIRED. Bare name exported.
                                 # ^[a-z0-9]([a-z0-9_-]*[a-z0-9])?$
      path: tools/kubectl.tool.yaml   # REQUIRED. Package-root-relative,
                                       # §Secure Path Resolution-safe
dependencies: []                # OPTIONAL. MUST be absent or empty in the
                                 # MVP; non-empty is PKG-016
\end{minted}

Rules:
\begin{itemize}
  \item The package root is the directory containing \texttt{gert-package.yaml}.
  \item \texttt{exports.tools[].id} is the \emph{only} way a tool becomes
        visible. A \texttt{.tool.yaml} inside the package that is not
        exported is invisible (\texttt{PKG-011} if referenced). There is
        \textbf{no} implicit \texttt{tools/} scanning inside a package ---
        implicit scanning is what turns package contents into an accidental
        API surface.
  \item The referenced \texttt{.tool.yaml}'s \texttt{meta.name} (or
        \texttt{name}) \textbf{MUST} equal the export \texttt{id}
        (\texttt{PKG-023}).
  \item The qualified catalog name is \texttt{<meta.name>/<export id>}.
  \item Non-tool assets are not describable in the MVP. Any other key under
        \texttt{exports} is a schema error.
\end{itemize}

\subsection{Project binding --- \texttt{.gert/config.yaml}}
\label{sec:tool-project-config}

Schema: \texttt{design/gert/schemas/project-config.v1.schema.json}
(\texttt{apiVersion: config/v1}). \texttt{.gert/config.yaml} previously had
no normative schema; it gets one here, covering only the fields this
specification touches. Other existing keys (\texttt{defaults},
\texttt{tools.mcp-servers}) are carried over as free-form objects.

\begin{minted}{yaml}
apiVersion: config/v1
requires:
  - package: acme.incident-tools   # REQUIRED
    version: "^1.4.0"              # REQUIRED (§Version Constraint Grammar)
    path: ./vendor/acme-incident-tools   # REQUIRED in the MVP; §Secure Path
                                          # Resolution-safe; MUST contain
                                          # gert-package.yaml
tool-paths:
  - ./ops/tools
\end{minted}

\texttt{path} is \textbf{REQUIRED} in the MVP precisely because there is no
remote catalog. When catalogs land post-MVP, \texttt{path} becomes optional
and its absence will mean ``resolve from catalog'' --- that is the designed
extension point, and it is forward-compatible without a schema break.

\subsection{Runbook binding --- additions to \texttt{runbook.v1.schema.json}}
\label{sec:tool-runbook-binding}

\begin{minted}{yaml}
requires:                    # NEW top-level, optional, array
  - package: acme.incident-tools
    version: "^1.4.0"
    path: ./vendor/acme-incident-tools   # optional here; if omitted, MUST
                                          # be supplied by project config
toolRefs:
  - name: kubectl
    package: acme.incident-tools
    version: "^1.4.0"
\end{minted}

See \S\ref{sec:tool-bindings} for the full disposition of \texttt{requires:}
and \texttt{toolRefs:}, including the duplicate/conflict rules
(\texttt{PKG-001}, \texttt{PKG-002}, \texttt{PKG-005}, \texttt{PKG-024}).

\subsection{Package map / lock --- \texttt{.gert/packages.lock.json}}
\label{sec:tool-package-lock}

Schema: \texttt{design/gert/schemas/package-lock.v1.schema.json}
(\texttt{apiVersion: package-lock/v1}). Generated, committed, and verified.
This is \textbf{not} a solver output --- there is nothing to solve in the
MVP --- it is a \textbf{frozen integrity record}.

\begin{minted}{json}
{
  "apiVersion": "package-lock/v1",
  "generatedBy": "gert/<version>",
  "packages": [
    {
      "name": "acme.incident-tools",
      "version": "1.4.2",
      "root": "vendor/acme-incident-tools",
      "digest": "sha256:<64 hex>",
      "exports": [
        { "id": "kubectl", "path": "tools/kubectl.tool.yaml",
          "digest": "sha256:<64 hex>" }
      ]
    }
  ],
  "catalogDigest": "sha256:<64 hex>"
}
\end{minted}

\begin{itemize}
  \item \texttt{root} is workspace-relative POSIX. Packages outside the
        workspace are recorded as
        \texttt{"external:sha256:<digest of the absolute realpath, UTF-8,
        NFC>"} --- never as a raw absolute path, because absolute paths leak
        usernames and host layout into committed artefacts and traces.
  \item Lock presence is \textbf{OPTIONAL} for \texttt{gert run}. When
        present, a mismatch is \texttt{PKG-009}.
  \item \texttt{gert verify} (\S\ref{sec:supply-chain}) is extended to check
        the lock.
\end{itemize}

% ═════════════════════════════════════════════════════════════════════════
\section{Action Substitution}
\label{sec:tool-substitution}
% ═════════════════════════════════════════════════════════════════════════

A tool action may be implemented by a GERT runbook instead of a process ---
this is called \emph{substitution}.

\subsection{Declaration}

Substitution is declared \textbf{in the \texttt{.tool.yaml}}, per action ---
never in the calling runbook. The caller \textbf{MUST NOT} be able to see or
change how an action is implemented.

The per-action key is \texttt{execute:}, \textbf{not} \texttt{impl:}. This
tool definition already has a top-level \texttt{impl:} block
(\S\ref{sec:tool-definition-schema}) for per-platform mobile dispatch
(\texttt{impl.ios}, \texttt{impl.android}, \texttt{transport: native-sdk},
\texttt{handler:}). Reusing that key at the action level for an unrelated
concept --- which implementation kind backs a single action --- would give
one key two meanings in the same document with no way to tell them apart
lexically. \texttt{execute:} is the disambiguated, schema-compatible name
for the per-action concept:

\begin{minted}{yaml}
actions:
  - name: drain-node
    execute:
      kind: runbook                    # process | runbook  (default: process)
      path: runbooks/drain-node.yaml   # relative to the DIRECTORY OF THIS
                                        # .tool.yaml
    args:
      node: { type: string, required: true }
    outputs:
      drained: { type: boolean }
      # A string output MAY carry enum, e.g.:
      # drain_state: { type: string, enum: ["drained", "cordoned"] }
\end{minted}

\texttt{execute.path} resolves relative to the declaring \texttt{.tool.yaml}'s
directory --- \textbf{not} the calling runbook, \textbf{not} the workspace
root, \textbf{not} the process cwd. This is the ``runbook-relative''
contract: a package is relocatable because every internal reference is
package-internal. Resolution obeys \S\ref{sec:tool-secure-paths}.

\paragraph{Interaction with tool-level \texttt{impl:} (mobile dispatch).}
The two keys operate at different scopes and are independent: tool-level
\texttt{impl:} selects a transport for the \textbf{entire tool} on a given
mobile platform, before any individual action's \texttt{execute:} is
consulted. When a tool declares \texttt{impl.ios} or \texttt{impl.android},
invocation on that platform \textbf{always} dispatches to the registered
native handler for \textbf{every} action of that tool, and
\texttt{execute.kind: runbook} on any of that tool's actions is
\textbf{ignored on that platform} --- substitution is a non-mobile,
non-native-SDK execution path. A tool that declares both a mobile
\texttt{impl:} block and a \texttt{kind: runbook} action is valid: the
substitute is used whenever the invoking host is not one of the platforms
listed under \texttt{impl:} (i.e. on desktop/server hosts), and the native
handler is used on the declared mobile platforms. This is stated normatively
here because the two blocks previously shared the name \texttt{impl:} with
no cross-reference, leaving this interaction undefined.

\subsection{Input contract}

\begin{itemize}
  \item The substitute runbook's \texttt{inputs:} map \textbf{MUST} be a
        \textbf{superset-free exact match by name} with the action's
        \texttt{args:} (\S\ref{sec:tool-definition-schema}): same key set, no
        extras, no omissions (\texttt{PKG-013}).
  \item Types \textbf{MUST} be identical PJVM types. No widening, no
        coercion (\texttt{PKG-013}).
  \item If either side declares \texttt{enum} on a shared \texttt{args.<name>},
        the other side \textbf{MUST} also declare \texttt{enum} with the
        \textbf{same member set} (order-insensitive set equality,
        \texttt{design/gert/sections/03-schema-vnext.tex} \S\ref{para:enum-constraint}).
        Neither side declaring \texttt{enum} is OK. One side declaring it
        while the other does not, or both declaring unequal sets, is
        \texttt{PKG-013} --- \textbf{no variance}: a narrower substitute
        would reject values the caller's published contract permits, and a
        wider substitute would silently accept values the published
        contract forbids and the caller's own plan-time \texttt{ENUM-007}
        check would have rejected as a literal. Both directions are
        contract breaks, not just one.
  \item \texttt{required} on the action side \textbf{MUST} imply
        \texttt{required} on the substitute side.
  \item Substitute inputs \textbf{MUST} have \texttt{from:} absent or
        \texttt{context}. \texttt{from: prompt} or \texttt{from: env} in a
        substitute is \texttt{PKG-026} --- a tool action must not covertly
        prompt the operator or read host environment.
  \item Argument values are bound positionally-by-name from
        \texttt{step.tool.args} after GIS interpolation in the
        \textbf{caller's} scope, exactly as for a process action
        (\S\ref{sec:tool-definition-schema} states that \texttt{args:} is
        the single vocabulary for \texttt{step.tool.args} regardless of
        \texttt{execute.kind}). The substitute receives values, never
        expressions. The caller's variable scope is not visible inside the
        substitute.
\end{itemize}

\subsection{Output contract}
\label{sec:tool-substitution-outputs}

\begin{itemize}
  \item The substitute's own \texttt{outputs:} map (its runbook-level
        \texttt{outputs:}, resolved against its own execution context at
        completion) \textbf{MUST} exactly match the action's \texttt{outputs:}
        key set and types (\texttt{PKG-013}).
  \item If either side declares \texttt{enum} on a shared
        \texttt{outputs.<name>}, the same set-equality rule as the Input
        contract applies (\texttt{PKG-013}, no variance): both sides
        declare an equal member set, or neither declares one.
  \item Every declared output \textbf{MUST} be produced on every terminal
        path of the substitute, or declare a default. Statically
        unproducible outputs are \texttt{PKG-027}. If the declared default
        is on an \texttt{enum}-constrained output, it MUST itself be a
        member (\texttt{ENUM-006}); the two checks are independent and both
        fire where applicable.
  \item An \texttt{enum}-constrained output's produced value is checked
        against its member set at \textbf{production time}, when it
        resolves at the substitute's own completion --- before it crosses
        the caller boundary or reaches \texttt{outputs.<name>} capture at
        the call site. A non-member value is \texttt{ENUM-009}, and the
        output is never emitted with a warning in its place.
  \item Captures inside the substitute are \textbf{not} visible to the
        caller. Only declared outputs cross the boundary. The trace shows
        the substitute's steps as nested spans (\S\ref{sec:pkg-events}),
        so auditability is preserved without leaking scope.
\end{itemize}

\paragraph{Caller-visible capture namespace (the substituted-action output
surface).} A substituted action's declared \texttt{outputs:} cross the
caller boundary under a new, dedicated GCP local-capture root:
\texttt{outputs.<name>} (design/gert/grammar/gcp.ebnf \S3.4a Substituted
Tool Action Output Capture, \texttt{LocalOutputs}). This is a distinct
surface from the
process-action capture surface (\texttt{stdout}, \texttt{exit\_code},
\texttt{json.*}, \texttt{yaml.*}): a substituted action never produces
stdout, an exit code, or a raw byte stream, so those roots \textbf{MUST NOT}
be used in a \texttt{capture:} block attached to a \texttt{tool} step whose
resolved action is \texttt{execute.kind: runbook} --- doing so is
\texttt{GCP-PARSE-001} (unknown source for this step), because the value
simply does not exist at that path. Concretely, for the reference
\texttt{drain-node} action (whose action declares
\texttt{outputs: \{ drained: \{ type: boolean \} \} }), the caller writes:

\begin{minted}{yaml}
capture:
  node_drained: outputs.drained
\end{minted}

\texttt{outputs.<name>} resolves to exactly the value the substitute's own
\texttt{outputs.<name>.value} produced (after that runbook's own GIS
evaluation completes), typed per the action's \texttt{outputs:} declaration.
\texttt{<name>} \textbf{MUST} be one of the action's declared output names;
any other name under \texttt{outputs.*} is \texttt{GCP-RESOLVE-002} (path
segment not found). No GDP suffix beyond the output name is supported in the
MVP, since substituted outputs are declared as scalars
(\S\ref{sec:tool-substitution}); a future revision may extend
\texttt{outputs.<name>} with a GDP suffix for structured output types.
\texttt{outputs.<name>} is capturable only in the \emph{same} step's
\texttt{capture:} block: \texttt{StepCapture} (\S\ref{sec:tool-substitution-outputs},
design/gert/grammar/gcp.ebnf \S3.1) gains no \texttt{step.\{id\}.outputs.\{name\}}
form in this MVP, so a peer step cannot re-read a substituted action's
output by referencing the tool step's id; a caller that needs the value
elsewhere in the flow must re-capture it locally (e.g.\ into a \texttt{vars:}
entry) from the step that owns it. This is a deliberate MVP scope choice,
not an oversight.

\subsection{Replay semantics}
\label{sec:tool-substitution-replay}

Replay mode (\texttt{gert exec --mode replay --scenario <file>},
design/gert/sections/13-evidence-tracing-resumption.tex \S Scenario File
Format, \S Replay Semantics) executes the runbook against a recorded
scenario file rather than the live environment: \texttt{cli} steps and
process-backed \texttt{tool} steps have their underlying command
invocation intercepted and answered from the scenario's \texttt{commands:}
list (matched by \texttt{argv}) instead of forking a real subprocess, and
\texttt{invoke} steps execute their child runbook in replay mode against
that same scenario. Action substitution (\texttt{execute.kind: runbook})
is resolved the same way as \texttt{invoke} for this purpose, and is
\textbf{not} treated as an ordinary process-backed \texttt{tool} step:
the substitute runbook is executed in replay mode against the caller's
scenario, and each of its own \texttt{cli}/\texttt{tool} steps is matched
individually against that scenario's \texttt{commands:} list, exactly as
if the substitute had been reached via an \texttt{invoke} step. Governance
rules for the substitute (\S\ref{sec:tool-substitution-governance}) are
evaluated normally in replay; replay does not bypass governance. Likewise,
\textbf{replay does not bypass enum validation}
(\texttt{design/gert/sections/03-schema-vnext.tex} \S\ref{para:enum-constraint}):
any \texttt{args.<name>}/\texttt{outputs.<name>} value answered from the
scenario's \texttt{commands:} list is a bound value like any other and is
checked against its declared \texttt{enum} identically to a live run; a
scenario recording an off-enum value fails replay with the same
\texttt{ENUM-008}/\texttt{ENUM-009} the live path would have raised. There is
no single ``recorded response'' attached to the substituted action as a
whole, because the action itself has no process-shaped invocation of its
own to match --- only its nested steps do.

This is directly load-bearing on \S7.6's package-drift semantics: replay
determinism is defined by the scenario file, not by the on-disk package
catalog (AR-TP §7.6) --- the substitute's own package/tool digests may
have drifted since the scenario was recorded, and replay proceeds anyway
because it never re-resolves or re-reads the substitute from the current
catalog to decide what to execute; it only needs the substitute
runbook's \emph{structure} (which steps exist, in what order) to know
which scenario commands to match against, which is why a digest mismatch
in replay mode is a non-fatal \texttt{replay/packageDrift} warning rather
than a hard refusal (design/gert/sections/06-tool-runtime.tex \S Digest,
Evidence, Trace, Resume; design/gert/sections/13-evidence-tracing-resumption.tex
\S Package Resumption Contract). This is distinct from the trace-event
replay described in design/gert/sections/07-runtime-events.tex \S
Determinism and Replay Semantics, which re-emits a historical event
stream for adapter/UI rendering of a completed run and performs no step
execution of any kind; that mechanism is unaffected by substitution
because it never invokes anything, tool-backed or otherwise.

\subsection{Governance contract}
\label{sec:tool-substitution-governance}

Substitution \textbf{MUST NOT} be a privilege-escalation path. The
substitute's effective policy is \textbf{derived}, never adopted. See
\texttt{design/gert/sections/12-governance-policy.tex}
\S Substitution Governance Composition for the full normative composition
table (invariants: \texttt{require\_approval} OR, \texttt{deny\_commands}
union, \texttt{deny\_env\_vars} union, \texttt{allow\_commands}
\textbf{intersection}, \texttt{redact} union, caller rules evaluated before
substitute rules with no deny-to-allow conversion, and capability closure
never a superset).

Any substitute construct whose composition would \textbf{widen} the
caller's effective policy is \texttt{PKG-014}, raised at plan time by static
comparison --- not at runtime. Widening detectable only at runtime is a
design defect; this is why \texttt{allow\_commands} is intersected rather
than merged.

The tool's own \texttt{governance.allowed-environments} and
\texttt{requires-approval} continue to gate the action as for
process-backed actions. A \texttt{dry-run} run \textbf{MUST NOT} execute a
substitute's side-effecting steps.

\subsection{Action reachability}
\label{sec:tool-substitution-reachability}

At plan time, for every \texttt{step.tool} in every reachable step of every
runbook in the include graph:

\begin{enumerate}
  \item \texttt{tool.name} \textbf{MUST} be bound by a \texttt{toolRefs}
        entry in the \textbf{same runbook file} (\texttt{PLAN-010}).
  \item \texttt{tool.action} \textbf{MUST} be declared by the bound tool
        definition (\texttt{PKG-012}).
  \item If the action is \texttt{execute.kind: runbook}, the substitute
        \textbf{MUST} parse, validate, and satisfy the input/output/governance
        contracts above (\texttt{PKG-013}/\texttt{PKG-014}/\texttt{PKG-027}).
  \item Every name listed in \texttt{toolRefs[].actions} \textbf{MUST}
        exist (\texttt{PKG-012}), even if no step uses it.
\end{enumerate}

Reachability uses the existing static-reachability analysis behind
\texttt{PLAN-009}. Unreachable steps are still validated: validity is not
conditional on execution path.

\subsection{Recursion}

Substitution frames are tracked as \texttt{(package, tool, action)} triples.
A cycle is \texttt{PKG-015}, detected at plan time by DFS. The maximum
substitution depth in the MVP is \textbf{4}; exceeding it is
\texttt{PKG-028}. Depth is bounded because each frame multiplies the
governance-composition surface. Substitution depth is counted separately
from include depth, and both bounds apply independently
(\S\ref{sec:tool-packages-includes}).

% ═════════════════════════════════════════════════════════════════════════
\section{Secure Path Resolution}
\label{sec:tool-secure-paths}
% ═════════════════════════════════════════════════════════════════════════

These rules apply uniformly to \texttt{requires[].path},
\texttt{exports.tools[].path}, \texttt{execute.path}, \texttt{toolRefs[].path},
\texttt{tool-paths[]}, and include/import paths inside packages.

\subsection{Resolution base per path kind}
\label{sec:tool-path-bases}

Every path kind below is resolved as
\texttt{realpath(join(dirname(referencing\_file), p))}, checked for
containment inside \texttt{realpath(containmentRoot)}, with
\texttt{referencing\_file} and \texttt{containmentRoot} named per kind in
Table~\ref{tab:tool-path-bases}. ``Package-internal'' kinds are
containment-checked against the \textbf{package root} and a violation is
\texttt{PKG-007} (\S\ref{sec:tool-path-containment}); ``workspace-level''
kinds are checked against the \textbf{workspace root} only for the purpose
of the external-root report (\texttt{PKG-W003}, \S\ref{sec:tool-path-symlinks}) and MAY legitimately resolve outside it --- they are
never rejected merely for escaping the workspace, since they are operator
configuration, not package content.

\begin{center}\small
\begin{tabular}{p{3.4cm}p{4.6cm}p{3.9cm}p{1.6cm}}
\hline
\textbf{Path kind} & \textbf{\texttt{referencing\_file}} & \textbf{\texttt{containmentRoot}} & \textbf{Class} \\
\hline
\texttt{requires[].path} (project scope, \texttt{.gert/config.yaml}) &
  the workspace root itself (i.e.\ \texttt{p} is workspace-relative, not
  relative to \texttt{.gert/}) &
  workspace root &
  workspace-level \\
\texttt{requires[].path} (runbook scope) &
  the declaring runbook file &
  workspace root &
  workspace-level \\
\texttt{exports.tools[].path} &
  \texttt{gert-package.yaml} (i.e.\ the package root) &
  package root &
  package-internal \\
\texttt{execute.path} &
  the declaring \texttt{.tool.yaml} &
  package root (of the package that exports the declaring tool) &
  package-internal \\
\texttt{toolRefs[].path} (tier 4) &
  the declaring runbook file &
  workspace root &
  workspace-level \\
\texttt{tool-paths[]} &
  the workspace root itself &
  workspace root &
  workspace-level \\
Package-internal include/import &
  the including file &
  package root &
  package-internal \\
\hline
\end{tabular}
\captionof{table}{Resolution base and containment root per path kind.}
\label{tab:tool-path-bases}
\end{center}

Two consequences follow directly from this table, both required to resolve
the ambiguity in the reference example (R5) and in fixture \texttt{r23}:

\begin{itemize}
  \item \texttt{execute.path}'s \textbf{resolution base} (the declaring
        \texttt{.tool.yaml}'s own directory) is \textbf{not necessarily the
        same as its containment root} (the package root). A tool definition
        nested at \texttt{tools/subdir/foo.tool.yaml} may reference
        \texttt{execute.path: ../../runbooks/foo.yaml}: the resolution base
        is \texttt{tools/subdir/}, the resolved path is
        \texttt{<packageRoot>/runbooks/foo.yaml}, and containment is checked
        against \texttt{<packageRoot>}, which it satisfies. This is why the
        reference tool's \texttt{tools/kubectl.tool.yaml} may legitimately
        use \texttt{execute.path: ../runbooks/drain-node.yaml}: resolved
        against \texttt{tools/} (the file's own directory) it lands at
        \texttt{<packageRoot>/runbooks/drain-node.yaml}, which is inside the
        package root and therefore \textbf{not} \texttt{PKG-007}. Evaluating
        the same literal string against the package root directly (i.e.\
        treating \texttt{execute.path} as package-root-relative rather than
        declaring-file-relative) would incorrectly compute
        \texttt{<packageRoot>/../runbooks/drain-node.yaml} and reject it ---
        that computation uses the wrong resolution base, not the wrong
        containment root, and is not how this specification defines
        resolution.
  \item Runbook-scope \texttt{requires[].path} is declaring-runbook-file
        relative, exactly like fixture \texttt{r23}'s
        \texttt{path: ../../../schemas/examples/acme-incident-tools}
        (resolved against
        \texttt{design/gert/testdata/runbooks/r23-tool-package/}, three
        levels up, landing at
        \texttt{design/gert/schemas/examples/acme-incident-tools}). Project-scope
        \texttt{requires[].path} is workspace-root relative, matching the
        \texttt{./vendor/...} convention in \S\ref{sec:tool-project-config}.
        These are two different bases for the textually identical field
        name at two different declaration scopes, and are now stated
        explicitly rather than left to inference.
  \item Table~\ref{tab:tool-path-bases} classifies \texttt{toolRefs[].path}
        as workspace-level (containment root = workspace root) because an
        \textbf{ordinary, top-level} runbook may point it at any tier-4
        \texttt{.tool.yaml} on disk. When the \textbf{declaring runbook file
        is itself package-internal} (i.e.\ it is reached only via a
        package's \texttt{exports.tools[].path} or an \texttt{execute.path},
        never run directly) --- as with the reference substitute
        \texttt{runbooks/drain-node.yaml} --- the blanket package rule of
        \S\ref{sec:tool-packages-includes} (``packages MUST NOT contain
        runbooks that include files outside the package root'') overrides
        the general classification: that file's own
        \texttt{toolRefs[].path} and \texttt{requires[].path} entries are
        package-internal for containment purposes, checked against the
        package root, not the workspace root. This is why
        \texttt{drain-node.yaml} binds its own \texttt{kubectl} reference
        via \texttt{toolRefs[].path: ../tools/kubectl.tool.yaml} (a sibling
        export inside the same package) rather than a self-referential
        \texttt{requires:} entry naming its own package --- both are valid
        under this rule, and the direct-path form avoids the conceptually
        circular case of a package requiring itself.
\end{itemize}

\subsection{Syntax rules}

Checked \textbf{before} touching the filesystem:

\begin{enumerate}
  \item The separator in authored artefacts is \texttt{/}. A backslash in an
        authored path is \texttt{PKG-007} --- not normalised, because
        \texttt{a\textbackslash{}b} is a legal single filename on POSIX.
  \item Absolute paths are rejected in \textbf{package-internal} references
        (\texttt{exports}, \texttt{execute}, package-internal includes):
        \texttt{PKG-007}. Absolute paths are permitted in
        \texttt{requires[].path} and \texttt{tool-paths[]} (workspace-level
        operator configuration) but are then subject to \S\ref{sec:tool-path-containment}.
  \item Rejected everywhere: drive-relative (\texttt{C:foo}), UNC
        (\texttt{\textbackslash{}\textbackslash{}host\textbackslash{}share}),
        extended-length (\texttt{\textbackslash{}\textbackslash{}?\textbackslash{}},
        \texttt{\textbackslash{}\textbackslash{}.\textbackslash{}}), device
        namespaces, \texttt{\textasciitilde{}} expansion, environment-variable
        expansion, URI schemes, NUL and control characters, and any segment
        equal to \texttt{.} or \texttt{..} \textbf{after} normalisation that
        escapes the root.
  \item Rejected filenames (all platforms, for portability): reserved DOS
        names \texttt{CON PRN AUX NUL COM1--COM9 LPT1--LPT9} with or without
        extension, and any segment with a trailing dot or trailing space
        (\texttt{PKG-018}).
  \item Paths are NFC-normalised for comparison. Two exports whose paths
        differ only by Unicode normalisation form, or only by ASCII case,
        are \texttt{PKG-018} --- this makes a package that works on Linux
        but breaks on Windows/macOS fail at authoring time, on every
        platform.
\end{enumerate}

\subsection{Containment rule}
\label{sec:tool-path-containment}

For every package-internal reference: \texttt{realpath(join(referencing\_dir, p))}
\textbf{MUST} be lexically inside \texttt{realpath(containmentRoot)}, where
\texttt{referencing\_dir} and \texttt{containmentRoot} are the values named
per kind in Table~\ref{tab:tool-path-bases}
(\S\ref{sec:tool-path-bases}), compared segment-wise on NFC-normalised,
case-folded segments when the host filesystem is case-insensitive, and
exactly otherwise. Violation is \texttt{PKG-007}.

Containment is checked on the \textbf{resolved} path, after link
resolution --- checking the lexical path only is the classic bypass.

For \textbf{workspace-level} kinds (\texttt{requires[].path} at either
scope, \texttt{toolRefs[].path} in an ordinary top-level runbook, and
\texttt{tool-paths[]}) the same computation is performed against the
workspace root, but the outcome is a \textbf{report, not a rejection}: a
resolved path outside the workspace root is recorded as an external root
and warned as \texttt{PKG-W003} (\S\ref{sec:tool-path-symlinks}, rule 6).
\texttt{PKG-007} \textbf{MUST NOT} be raised for a workspace-root escape by
any workspace-level kind; for those kinds \texttt{PKG-007} remains reachable
only via the syntax rules above. If the escaping target does not contain a
resolvable package, the resulting error is \texttt{PKG-001}, not
\texttt{PKG-007}.

\subsection{Symlinks, junctions, reparse points}
\label{sec:tool-path-symlinks}

\begin{enumerate}
  \item Symlinks are \textbf{resolved, then containment-checked}. A link
        whose final target is inside the root is allowed. A link whose
        final target is outside the root is \texttt{PKG-008} --- an
        \textbf{error}, never a silent skip, because silent skipping
        produces a catalog that differs between machines.
  \item Windows directory junctions, mount points, and any
        \texttt{IO\_REPARSE\_TAG\_*} reparse point are treated exactly as
        symlinks. AppExecLink and OneDrive placeholder reparse points are
        \texttt{PKG-008}.
  \item Link chains are bounded at \textbf{8} hops; exceeding it is
        \texttt{PKG-008}. Cycles are \texttt{PKG-008}.
  \item Resolution \textbf{MUST} be performed with a TOCTOU-resistant
        sequence: resolve once, then read via the resolved handle/path; the
        digest (\S\ref{sec:tool-package-digest}) is computed over the same
        resolved bytes used for loading.
  \item Hard links are not detectable and are out of scope; the digest
        covers content, which is the property that matters.
  \item \texttt{requires[].path} and \texttt{tool-paths[]} that resolve
        outside the workspace are permitted but \textbf{MUST} be recorded in
        the lock and trace in the \texttt{external:sha256:...} form
        (\S\ref{sec:tool-package-lock}), and \texttt{gert verify}
        \textbf{MUST} report them as a distinct ``external root'' finding
        (\texttt{PKG-W003}).
\end{enumerate}

% ═════════════════════════════════════════════════════════════════════════
\section{Digest, Evidence, Trace, Resume}
\label{sec:tool-package-digest}
% ═════════════════════════════════════════════════════════════════════════

\subsection{File digest}

\texttt{sha256} over the file's \textbf{raw bytes}. No line-ending
normalisation, no whitespace normalisation, no YAML canonicalisation.
Encoded as \texttt{sha256:} followed by 64 lowercase hex characters.

Any normalisation is a second parser, and a second parser is a second bug.
The Windows CRLF hazard is handled by a documented \texttt{.gitattributes}
recommendation (\texttt{*.tool.yaml text eol=lf}, \texttt{gert-package.yaml
text eol=lf}), not by a normalisation rule.

\subsection{Package digest (deterministic, order-free)}

Given package root $R$, the digest set is: \texttt{gert-package.yaml}, every
\texttt{exports.tools[].path}, and every file transitively referenced by an
exported tool via \texttt{execute.path} and package-internal includes. Files
present in the package but not in this closure are \textbf{not} hashed ---
the digest covers the API surface, not incidental contents.

\begin{minted}{text}
for each file f in closure:
    line(f) = hex(sha256(bytes(f))) + "  " + relpath_posix_nfc(f, R)
sort lines ascending by UTF-8 byte order of the whole line
blob = concat(line + "\n" for each sorted line)
packageDigest = "sha256:" + hex(sha256(blob))
\end{minted}

Two spaces separate the digest and the path (the coreutils
\texttt{sha256sum} convention), so the blob is independently reproducible
with standard tools.

\subsection{Catalog digest}

Same construction over lines of
\texttt{qualifiedName + "  " + packageDigestOrFileDigest + "  " + tierNumber},
sorted ascending. \texttt{catalogDigest} is the single value that proves two
runs saw the same tool universe.

\subsection{Trace, evidence, and resume}

New trace events (\texttt{package/resolved}, \texttt{catalog/frozen},
\texttt{tool/substituted}), the resume refusal contract, and the replay
non-fatal drift contract are specified normatively in
\texttt{design/gert/sections/07-runtime-events.tex}
\S Package and Substitution Events and
\texttt{design/gert/sections/13-evidence-tracing-resumption.tex}
\S Package Resumption Contract. In summary:

\begin{itemize}
  \item Package roots, digests, and the catalog digest are part of the run
        manifest and are covered by the existing HMAC trace chaining
        (\S\ref{sec:supply-chain}). Package paths are subject to the
        standard redaction pipeline; the \texttt{external:sha256:} form is
        applied \emph{before} redaction, not after.
  \item On \texttt{gert exec --resume}, a package or catalog digest mismatch
        is a hard refusal (\texttt{PKG-009}); \texttt{--allow-package-drift}
        overrides and emits an auditable
        \texttt{governance/packageDriftAccepted} event.
  \item In replay mode, a digest mismatch is non-fatal: it emits
        \texttt{replay/packageDrift} and replay proceeds, because replay
        determinism is defined by the scenario file, not by the on-disk
        catalog.
\end{itemize}

% ═════════════════════════════════════════════════════════════════════════
\section{Includes, Lexical Scoping, and Global Package Set}
\label{sec:tool-packages-includes}
% ═════════════════════════════════════════════════════════════════════════

\textbf{Tool name binding is lexically scoped to the runbook file that
declares it.} An included child runbook does \textbf{NOT} inherit the
parent's \texttt{toolRefs}. If a child uses \texttt{kubectl}, the child
declares \texttt{kubectl} in its own \texttt{toolRefs}. This is a decisive,
intentional call: dynamic (inherited) scoping would make a child runbook's
meaning depend on who included it, which destroys both composability and
static reachability analysis. See
\texttt{design/gert/sections/02-architecture.tex}
\S Lexical Tool-Name Scoping for the executor-level statement of this rule.

Packages are \textbf{globally resolved, per run, at the root runbook}. The
effective package set is the union of the project-level \texttt{requires:}
and the runbook-level \texttt{requires:} of \emph{every} runbook in the
include/import closure, computed at plan time. Constraints for the same
package name across the closure are intersected
(\S\ref{sec:tool-bindings}); an empty intersection is \texttt{PKG-002}.
Consequently, a package is loaded once per run, at one version, with one
digest --- two runbooks in one run cannot use two versions of the same
package. This is a deliberate MVP restriction; multi-version coexistence
requires a tier-aware qualified namespace and is deferred
(\S\ref{sec:tool-packages-nongoals}).

The include closure \textbf{MUST} be fully enumerable at plan time. Lazy
includes (\texttt{expand: lazy}) are still \emph{resolved and
package-analysed} at plan time even though they are \emph{expanded} at run
time. A lazily-expanded include that would introduce a package not in the
frozen set is \texttt{PKG-017}. Packages \textbf{MUST NOT} contain runbooks
that include files outside the package root (\S\ref{sec:tool-path-containment});
a package's runbooks are self-contained. \texttt{imports:} (alias$\to$path)
remains a runbook-local alias map for child runbooks and is unrelated to
packages; it \textbf{MUST NOT} be repurposed for package aliasing.
Alias-based package naming is out of scope for the MVP. Maximum include
depth is unchanged; substitution depth (\S\ref{sec:tool-substitution}) is
counted separately and both bounds apply independently.

% ═════════════════════════════════════════════════════════════════════════
\section{Non-Goals}
\label{sec:tool-packages-nongoals}
% ═════════════════════════════════════════════════════════════════════════

The following are explicit non-goals of this MVP. GERT is pre-1.0 with no
external consumers (ratified OQ-M4, ``just ship''); the MVP pins determinism
and safety, not distribution.

\begin{itemize}
  \item \textbf{Remote catalogs and registries.} All packages are resolved
        from a local, workspace-relative or absolute filesystem path
        (\S\ref{sec:tool-project-config}). There is no network fetch, no
        registry protocol, and no package index.
  \item \textbf{Publishing / packaging workflows.} There is no
        \texttt{gert publish}, no package build step, and no artefact
        upload path.
  \item \textbf{Transitive package dependencies.} A package's
        \texttt{dependencies:} field \textbf{MUST} be absent or empty
        (\texttt{PKG-016}). Consumers flatten any indirect tool dependency
        into their own \texttt{requires:} list.
  \item \textbf{Tool aliases.} \texttt{toolRefs[].alias} is removed
        (\texttt{PKG-021}); \texttt{imports:} alias maps are unrelated to
        packages and are not repurposed for them.
  \item \textbf{Action allowlists as an enforcement mechanism.}
        \texttt{toolRefs[].actions} is documentation and a plan-time
        reachability assertion only (\texttt{PKG-W001}); it does not gate
        invocation.
  \item \textbf{Packaging of non-tool assets.} Docs, dashboards, policy
        bundles, and binaries-as-payload are not describable by
        \texttt{gert-package.yaml} in the MVP.
  \item \textbf{Multi-version coexistence.} A single run resolves at most
        one version of a given package (\S\ref{sec:tool-packages-includes}).
  \item \textbf{Enum member sets in the package lock or any digest beyond
        file bytes.} A ``mock package'' is not a spec concept --- it is an
        ordinary package bound via \texttt{.gert/config.yaml}
        \texttt{requires[].path} or a \texttt{--package-map} override.
        \texttt{enum} validation is binding-independent and applies
        identically regardless of which binding is resolved
        (\S\ref{sec:tool-substitution}); a mock declaring a different
        \texttt{enum} set than the real package for the same action is an
        authoring error enforced by a required conformance vector class
        (\texttt{ENUM-MOCK}, \texttt{design/gert/conformance/tv-enum.yaml}),
        \textbf{not} by a runtime signature comparison, a package lock
        field, or new digest machinery: an enum edit inside a package
        already changes the export digest and is already \texttt{PKG-009}.
\end{itemize}

\subsection{Compatibility: No Shims, One Intentional Breaking Change}
\label{sec:tool-packages-compat}

Per AR-TP-10, GERT is pre-1.0 and this specification introduces \textbf{no
compatibility shims}: there is no dual-read of \texttt{toolPackages:} and
\texttt{requires:}, no silent \texttt{alias} tolerance, and no legacy
resolution-order fallback path. \texttt{source} and \texttt{actions} remain
schema-valid only per the deprecation schedule in
\S\ref{sec:tool-bindings}; they carry no special-case resolution behaviour.

Exactly one behaviour change in this specification is intentionally
breaking rather than additive: \textbf{undeclared cross-tier shadowing,
which previously resolved silently by ``first match wins'' /
``later entries win'' precedence, is now \texttt{PKG-022}, a hard plan-time
error} (\S\ref{sec:tool-precedence}). A runbook that previously relied on
one tier's tool silently shadowing another's same-named tool without an
explicit \texttt{toolRefs:} disambiguation will now fail to plan. This is a
deliberate correctness fix, not an oversight: silent shadowing is exactly
the class of non-determinism this MVP exists to eliminate, and pre-1.0 is
the only opportunity to remove it without a deprecation window.

\section{Tool Versioning}
\label{sec:tool-versioning}

Tool definitions carry a \texttt{meta.version} field, validated as strict
SemVer 2.0.0 (\S\ref{sec:tool-version-grammar}). When multiple tool
definitions could in principle share the same name, the catalog's tier and
collision rules (\S\ref{sec:tool-tiers}, \S\ref{sec:tool-precedence}) ---
not version comparison --- determine which definition is bound to that
name; there is no ``highest version wins'' fallback. Explicit version
constraints are declared via \texttt{requires:}
(\S\ref{sec:tool-bindings}) and, optionally, per tool reference via
\texttt{toolRefs[].version} (\S\ref{sec:tool-runbook-binding}):

\begin{minted}{yaml}
requires:
  - package: acme.incident-tools
    version: "^1.4.0"
\end{minted}
